\begin{center}
    {\Large\itshape Alban en fin de siècle (XV\textsuperscript{ème})}
\end{center}

\bigskip

En cette fin de de XV\textsuperscript{ème} siècle, la paix semble donc bien assise. Aussi, en 1483 (nous sommes sous le règne de Charles VIII), Jean d'Armagnac, qui avait dû s'exiler tandis que son frère avait porté son cou sur le billot, veut réintégrer son évêché. Mais cette entrée dans la ville épiscopale de Castres est l'occasion de troubles. Partout, on se soulève contre celui qui a été imposé par le gouvernement royal comme gouverneur de Castres : Foffille de \footnote{nom illisible}. Et notre cité n'est pas la dernière à lever l'étendard de la révolte. La barre recommence. Ces « Albans » ont du sang de Vercingétorix dans les veines. César ne disait-il pas que si les Germains combattaient pour la rapine (ad praedam), les Gaulois, eux, c'était pour la liberté (pro libertate). Pauvres « Albans » de ce temps-là. L'affaire tourna mal pour eux. Bollille de Juge, après s'être emparé de Viane, prit d'assaut la petite place forte d'Alban. La cité rebelle fut pillée consciencieusement. Même son église ne fut pas épargnée. C'était en l'an 1484.

À part cet incident malheureux, il ne reste pas grand chose à glaner de notre histoire pour satisfaire notre curiosité. Citons cependant une délimitation entre les territoires de la paroisse N.D. d'Alban et ceux de St Michel de la Besse (actuellement la ferme St Michel, non loin du rocher de ce nom). Figurent dans l'acte passé devant Me Laurent, notaire : le représentant du chapitre de la cathédrale Sainte Cécile d'Albi à savoir le nouveau prieur de Négremont, pour le monastère St Benoît de Castres, son prieur, noble Raymond qu'accompagne Pierre ?alvan\footnote{Illisible} qui est de Castres. La limite fixée du côté du tènement appelé « lou Puech de la Bouyssière » (ne serait-ce pas ce qu'on appelle au roc des Anglais : la Puyssière ?...), confrontait vers le levant le château de Nougaret, situé au dessus du chemin d'Alban à Saint-André. Cette partie devait appartenir au prieur de Saint-Benoît, donc à l'abbaye, et la partie opposée située dans la juridiction d'Ambialet devait appartenir au chapitre cathédral de Ste Cécile d'Albi.
